\section{Comparison on a frozen ledger}\label{sec:frozen-frontier}

Section~\ref{sec:fixed-package} showed that growth requires switching packages. Different switches lead to different extensions, so the next question is how to compare extensions. We fix a ledger $L=(Y,K)$, a comparison domain $D$ and a cone $C$, and ask what efficiency under $L$ can and cannot tell us.

\subsection{Efficiency belongs to agreement classes}

\begin{proposition}[Dominance is invariant under cone agreement]\label{prop:dominance-congruence}
Let $L$ be invariant on $C$, and let $f\agree{C}f'$. Then for every operator $g$, $g\Dom_L f$ if and only if $g\Dom_L f'$.
\end{proposition}

\begin{proof}
Invariance gives $Y(f)=Y(f')$ and $K(f)=K(f')$. Dominance of $f$ by $g$ depends on $f$ only through these two numbers.
\leanref{\lid{frontier\_dominates\_congr\_right}}
\end{proof}

\begin{theorem}[Efficiency transfer]\label{thm:efficiency-transfer}
Let $L$ be invariant on $C$. If $\Eff_{L,D}(e)$, $e\agree{C}e'$ and $e'\in D$, then $\Eff_{L,D}(e')$.
\end{theorem}

\begin{proof}
Suppose some $g\in D$ dominates $e'$. By Proposition~\ref{prop:dominance-congruence}, applied to $e'\agree{C}e$, the same $g$ dominates $e$, contradicting efficiency of $e$.
\leanref{\lid{frontier\_efficiency\_transfer\_to\_cone\_equivalent}}
\end{proof}

Under an invariant ledger, then, efficiency is a property of agreement classes. The ledger cannot reward one representative of a class over another. This yields the basic replacement step used later.

\begin{lemma}[Restricted alignment]\label{lem:restricted-alignment}
Let $L$ be invariant on $C$, let $\mathcal{K}\subseteq D$, and let $e$ be efficient in $D$ and aligned with $\mathcal{K}$ on $C$. Then there is $c\in\mathcal{K}$ that is efficient in $D$ and satisfies $e\agree{C}c$.
\end{lemma}

\begin{proof}
Choose $c\in\mathcal{K}$ with $e\agree{C}c$. Since $c\in D$, Theorem~\ref{thm:efficiency-transfer} applies.
\leanref{\lid{restricted\_in\_house\_alignment\_lemma}}
\end{proof}

The conclusion keeps the agreement $e\agree{C}c$ with the original operator. Without it the statement would be trivial, since any efficient canonical operator is aligned with itself.

\subsection{What efficiency cannot do}

Lemma~\ref{lem:restricted-alignment} \emph{assumes} alignment. A natural hope is that efficiency could \emph{supply} it, so that being efficient under every reasonable ledger forces an operator into the class of some canonical operator. The next theorem shows that, without restrictions on the ledger, this happens only when the canonical family already meets every agreement class.

\begin{theorem}[Ledger boundary]\label{thm:ledger-boundary}
Let $\mathcal{K}\subseteq D$. The following are equivalent.
\begin{enumerate}
\item[\textup{(i)}] For every ledger $L$ invariant on $C$ and every $e$ efficient in $D$ under $L$, some $c\in\mathcal{K}$ is efficient in $D$ under $L$ and satisfies $e\agree{C}c$.
\item[\textup{(ii)}] Every $e\in D$ is aligned with $\mathcal{K}$ on $C$.
\end{enumerate}
\end{theorem}

\begin{proof}
(ii)$\Rightarrow$(i) is Lemma~\ref{lem:restricted-alignment}, because efficient operators lie in $D$. For (i)$\Rightarrow$(ii), fix $e\in D$ and use the \emph{indicator ledger} of its class: $Y(f)=1$ if $f\agree{C}e$ and $Y(f)=0$ otherwise, with $K\equiv 0$. Since $\agree{C}$ is an equivalence relation, this ledger is invariant on $C$. Nothing has yield above $1$, so $e$ is efficient in $D$. Statement~(i) now supplies $c\in\mathcal{K}$ with $e\agree{C}c$.
\leanref{\lid{universal\_ledger\_canonicality\_iff\_coverage}}
\end{proof}

\begin{corollary}[Counterledger]\label{cor:counterledger}
If some $e\in D$ is not aligned with $\mathcal{K}$ on $C$, there is a ledger invariant on $C$ under which $e$ is efficient in $D$ and has no canonical representative on~$C$.
\leanref{\lid{uncovered\_class\_has\_counterledger}}
\end{corollary}

Two remarks delimit the theorem. First, coverage does not require $\mathcal{K}=D$. A proper subfamily suffices as long as it meets every class, and the Lean regression tests contain such an example. Second, restricting the ledgers, for instance to computable ones, does not weaken the theorem. The direction (i)$\Rightarrow$(ii) needs only a single ledger, and the constant ledger $Y\equiv K\equiv 0$ already works. It is computable and invariant on every cone, and under it no operator dominates another, so every member of $D$ is efficient. Hence the equivalence holds for every class of ledgers that contains a constant ledger. The Lean proof uses the indicator ledger described above; this one-line variant is not separately formalized. The undecidability of cone agreement plays no role.

In short, efficiency can select canonical representatives beyond what coverage already provides only if the admissible ledgers exclude the constant ones, for example by requiring ledgers to reward strictly stronger extensions. Otherwise any route to canonical representatives must establish coverage of the agreement classes directly. Section~\ref{sec:conditional-lift} takes this route as an explicit assumption. Section~\ref{sec:arithmetic-instance} establishes the analogous fact in a concrete arithmetic case.

\subsection{Finite frontiers}

For finite candidate lists a frontier always exists. If a nonempty finite list of candidates carries rational yield and cost scores, some candidate is not dominated by any other. The proof is by induction on the list and uses that strict Pareto dominance is irreflexive and transitive. The set of nondominated candidates does not depend on the order in which the list is enumerated. (Lean: \lid{finite\_nondominated\_exists}, \lid{finite\_frontier\_enumeration\_invariant}.)

Existence is all one gets, however. A computer search over all $422$ monotone inflationary self-maps of the chains with $2$ to $6$ elements and of the Boolean lattices of rank $1$ to $3$ evaluated these maps under several benchmark weightings and cost functions. Uniqueness of the efficient map fails already in the smallest models, the two-element chain and lattice. Stability of the frontier under reweighting of the benchmark fails on the three-element chain and the rank-$2$ Boolean lattice. Containment of the frontier in a natural ``step'' family on chains fails on the three-element chain. These are finite computations. The enumeration of the $422$ maps was checked independently, and $128$ of the scoring and frontier computations were repeated in exact rational arithmetic. They are consistent with Theorem~\ref{thm:ledger-boundary}: efficiency by itself does not pick out a canonical operator.
